% GENERATED FILE. Edit canonical lessons, then run npm run generate:books.
% canonical-source-hash: fnv1a64:ce3148b972492a40
% canonical-lessons: KA-C205-kilu, KA-C205-turi, KA-C205-sosu, KA-C205-nadu, KA-C205-lattisu

\chapter{To pluck, To grate, To strain, To knead, To roll out}
\label{ch:ka-to-pluck-to-grate-to-strain-to-knead-to-roll-out}
\hlchaptermodality{205}

\begin{chapteropening}
\textbf{By the end of this chapter:} \emph{I can say to pluck, to grate, to strain, to knead and to roll out.}

Complete the last lesson of chapter 205: I can say to pluck, to grate, to strain, to knead and to roll out.
\end{chapteropening}

\section[\texorpdfstring{\kn{ಕೀಳು} (kīḷu)}{ಕೀಳು (kīḷu)}]{\kn{ಕೀಳು} (kīḷu) — to pluck, to pull out}

\label{lesson:KA-C205-kilu}



\begin{quote}
Before the new one: say the Kannada for to have fun, then the Kannada for to compete.
\end{quote}

\subsection*{You'll want to know: \kn{ಕೀಳು}}
\textbf{\kn{ಕೀಳು}} — \emph{kīḷu} — \textquotedblleft{}to pluck, to pull out\textquotedblright{}.

\begin{grammarlens}[title={What you've built}]
One more word for this chapter.
\end{grammarlens}

\subsection*{Your turn}
\begin{itemize}
\raggedright
  \item \emph{Say it:} \emph{kīḷu}
  \item \emph{Say it:} \emph{kīḷu}, once more
  \item \emph{Say it:} say \emph{kīḷu}
  \item \emph{Recall:} say the Kannada for to have fun, then the Kannada for to compete, then say \emph{kīḷu} again
\end{itemize}

\begin{tcolorbox}[breakable,colback=teal!4,colframe=teal!35!black,title={Before you move on}]
What does \kn{ಕೀಳು} mean? (\textquotedblleft{}To pluck, to pull out\textquotedblright{}.) Say it once more.
\end{tcolorbox}
\section[\texorpdfstring{\kn{ತುರಿ} (turi)}{ತುರಿ (turi)}]{\kn{ತುರಿ} (turi) — to grate (coconut)}

\label{lesson:KA-C205-turi}



\begin{quote}
Before the new one: say the Kannada for to compete, then the Kannada for to pluck.
\end{quote}

\subsection*{You'll want to know: \kn{ತುರಿ}}
\textbf{\kn{ತುರಿ}} — \emph{turi} — \textquotedblleft{}to grate (coconut)\textquotedblright{}.

\begin{grammarlens}[title={What you've built}]
One more word for this chapter.
\end{grammarlens}

\subsection*{Your turn}
\begin{itemize}
\raggedright
  \item \emph{Say it:} \emph{turi}
  \item \emph{Say it:} \emph{turi}, once more
  \item \emph{Say it:} say \emph{turi}
  \item \emph{Recall:} say the Kannada for to compete, then the Kannada for to pluck, then say \emph{turi} again
\end{itemize}

\begin{tcolorbox}[breakable,colback=teal!4,colframe=teal!35!black,title={Before you move on}]
What does \kn{ತುರಿ} mean? (\textquotedblleft{}To grate (coconut)\textquotedblright{}.) Say it once more.
\end{tcolorbox}
\section[\texorpdfstring{\kn{ಸೋಸು} (sōsu)}{ಸೋಸು (sōsu)}]{\kn{ಸೋಸು} (sōsu) — to strain, to filter}

\label{lesson:KA-C205-sosu}



\begin{quote}
Before the new one: say the Kannada for to pluck, then the Kannada for to grate.
\end{quote}

\subsection*{You'll want to know: \kn{ಸೋಸು}}
\textbf{\kn{ಸೋಸು}} — \emph{sōsu} — \textquotedblleft{}to strain, to filter\textquotedblright{}.

\begin{grammarlens}[title={What you've built}]
One more word for this chapter.
\end{grammarlens}

\subsection*{Your turn}
\begin{itemize}
\raggedright
  \item \emph{Say it:} \emph{sōsu}
  \item \emph{Say it:} \emph{sōsu}, once more
  \item \emph{Say it:} say \emph{sōsu}
  \item \emph{Recall:} say the Kannada for to pluck, then the Kannada for to grate, then say \emph{sōsu} again
\end{itemize}

\begin{tcolorbox}[breakable,colback=teal!4,colframe=teal!35!black,title={Before you move on}]
What does \kn{ಸೋಸು} mean? (\textquotedblleft{}To strain, to filter\textquotedblright{}.) Say it once more.
\end{tcolorbox}
\section[\texorpdfstring{\kn{ನಾದು} (nādu)}{ನಾದು (nādu)}]{\kn{ನಾದು} (nādu) — to knead (dough)}

\label{lesson:KA-C205-nadu}



\begin{quote}
Before the new one: say the Kannada for to grate, then the Kannada for to strain.
\end{quote}

\subsection*{You'll want to know: \kn{ನಾದು}}
\textbf{\kn{ನಾದು}} — \emph{nādu} — \textquotedblleft{}to knead (dough)\textquotedblright{}.

\begin{grammarlens}[title={What you've built}]
One more word for this chapter.
\end{grammarlens}

\subsection*{Your turn}
\begin{itemize}
\raggedright
  \item \emph{Say it:} \emph{nādu}
  \item \emph{Say it:} \emph{nādu}, once more
  \item \emph{Say it:} say \emph{nādu}
  \item \emph{Recall:} say the Kannada for to grate, then the Kannada for to strain, then say \emph{nādu} again
\end{itemize}

\begin{tcolorbox}[breakable,colback=teal!4,colframe=teal!35!black,title={Before you move on}]
What does \kn{ನಾದು} mean? (\textquotedblleft{}To knead (dough)\textquotedblright{}.) Say it once more.
\end{tcolorbox}
\section[\texorpdfstring{\kn{ಲಟ್ಟಿಸು} (laṭṭisu)}{ಲಟ್ಟಿಸು (laṭṭisu)}]{\kn{ಲಟ್ಟಿಸು} (laṭṭisu) — to roll out (dough)}

\label{lesson:KA-C205-lattisu}



\begin{quote}
Before the new one: say the Kannada for to strain, then the Kannada for to knead.
\end{quote}

\subsection*{You'll want to know: \kn{ಲಟ್ಟಿಸು}}
\textbf{\kn{ಲಟ್ಟಿಸು}} — \emph{laṭṭisu} — \textquotedblleft{}to roll out (dough)\textquotedblright{}.

\begin{grammarlens}[title={What you've built}]
One more word for this chapter.
\end{grammarlens}

\subsection*{Your turn}
\begin{itemize}
\raggedright
  \item \emph{Say it:} \emph{laṭṭisu}
  \item \emph{Say it:} \emph{laṭṭisu}, once more
  \item \emph{Say it:} say \emph{laṭṭisu}
  \item \emph{Recall:} say the Kannada for to strain, then the Kannada for to knead, then say \emph{laṭṭisu} again
\end{itemize}

\begin{tcolorbox}[breakable,colback=teal!4,colframe=teal!35!black,title={Before you move on}]
What does \kn{ಲಟ್ಟಿಸು} mean? (\textquotedblleft{}To roll out (dough)\textquotedblright{}.) Say it once more.
\end{tcolorbox}
